\section{Discussion: Architectural Stability}

The transition from an ephemeral component lifecycle to a persistent shell model
represents a structural shift in UI architecture. Traditional component-based
frameworks repeatedly instantiate and reconcile intermediate structures,
leading to dynamic allocation patterns and reconciliation overhead. 
Domphy instead establishes a stable structural topology during the
reification phase and maintains it throughout the application lifecycle.

By representing UI elements as flat compositional objects with intrinsic
reactive bindings, Domphy separates structural definition from runtime
mutation. The structural graph is fixed after initialization, while
state changes propagate exclusively through pre-registered atomic bindings.
This separation enhances architectural determinism and reduces runtime
structural variability.

From a cognitive perspective, the flat composition model reduces indirection.
Developers interact with a single object literal rather than navigating
multi-layered provider hierarchies or nested component wrappers.
The absence of recursive re-render cycles simplifies mental modeling of
state-to-view synchronization.

From a systems perspective, architectural stability yields predictable
performance characteristics. Since updates are constrained to bounded
atomic slots, latency does not scale with global tree size. This property
is particularly relevant for high-frequency interaction environments,
including data visualization dashboards, animation-heavy interfaces,
and direct-manipulation design tools, where frame-rate stability is
more critical than generalized state-tree abstraction.

However, this model introduces a space–time trade-off. Persistent bindings
increase memory residency due to retained reactive listeners.
While modern hardware environments typically tolerate this overhead,
large-scale applications must consider memory–CPU balance when adopting
the CARM model.
